\section{Two-dimensional acyclic resolutions: independent audit}
\label{sec:E-audit}

In this section a bound-$\mu$ resolution of a simply connected space $W$
means a group $G$, an actual two-dimensional CW model $X=K(G,1)$,
and a map $t:X\to W$ inducing integral homology isomorphisms, with
every homomorphism from $G$ to a group of cardinality at most $\mu$
trivial. The trivial group is allowed. All groups and cell sets are small.
We use $\mathfrak c=2^{\aleph_0}$.

\paragraph{Statement E.1 (verified input for the cardinal argument).}
Let $P$ be a group and $w_1,\ldots,w_m\in P*F(t_1,\ldots,t_n)$.
If their exponent-sum rows have rank $m$ over $\mathbb Q$, then
\[
 P\longrightarrow (P*F(t_1,\ldots,t_n))/
                   \langle\!\langle w_1,\ldots,w_m\rangle\!\rangle
\]
is injective. This is precisely Theorem~1 of the OpenAI preprint
\emph{Nonsingular systems of equations over arbitrary groups}, dated
5 October 2026 \cite[Theorem~1 and Remark~5.1]{EOpenAI2026}.
The complete sixteen-page source was obtained from the
official repository and read; it imposes no restriction on $P$.
The source is a preprint, and no claim of formal verification is made here.

Here are the proof checks relevant to using that theorem. For unitary
variables of sizes $k_1\geq\cdots\geq k_r$, rational independence of
the signed block multiplicity rows supplies nested nonzero minors.
The incidence space parametrizing a unitary matrix and an $l$-plane
fixed by it has codimension $l^2$. Its Poincar\'e dual has leading
exterior term a nonzero multiple of the product of the first $l$
primitive cohomology generators of the unitary group. Restriction to
$X\mapsto X\oplus(-I)$ verifies this coefficient by a single transverse
intersection; shorter exterior terms cannot contribute to the
top-dimensional pairing. The pairing for the product of the unitary
variable groups is the product of the indicated nested minors, hence
nonzero. This establishes the required simultaneous fixed-space bound.

For the operator step, work in the matrix algebra over the group von
Neumann algebra, with its faithful, unnormalized trace $\tau$.
Writing $h(e^{2\pi it})=t-\lfloor t\rfloor$ and
$\ell(U)=\tau(h(U))$, differentiate
$\tau(f(Ue^{2\pi is h(V)}))$ for smooth $f$ on the circle.
Cyclicity of the trace gives the derivative
$\tau(f'(Ue^{2\pi is h(V)})h(V))$; it is at most $\tau(h(V))$
when $f'\leq1$. Smoothing the sawtooth from the right, including at
the cut $1$, gives $\ell(UV)\leq\ell(U)+\ell(V)$ by bounded
convergence. For a weighted cyclic shift of length $a$, the phases of
the $a$ roots give $(a-1)/2$ plus the phase of the product of its weights.
These facts remain valid for uncountable coefficient groups: every
square-summable coefficient vector has countable support, and the trace
is still faithful and normal.

For a planar band diagram, independence of the exponent rows forces
equal positive and negative disk counts for each relator type. If there
are $2K$ disks, $b$ bands and $q$ boundary components, then
$2-q=2K-b$. The fixed-space bound and subadditivity give the upper
phase estimate $3b/2-K$. Assign a small negative phase to each
nonexceptional boundary and the compensating positive phase to the
exceptional boundary. The weighted-cycle formula tends to the same
constant plus the nonnegative phase of its coefficient word. Faithfulness
therefore makes that word the identity. Finally, a finite normal-closure
expression for a coefficient kernel element gives a punctured disk;
disjoint level arcs of the variable-circle maps supply planar bands
on which the coefficient map is constant. Removing components from
the inside out proves coefficient injectivity. This checks the source's
comparison from a kernel relation to its planar obstruction, rather
than assuming that an arbitrary pairing of letters is planar.

\paragraph{Statement E.2 (confirmation of INDLOCAL A.19).}
Let $X=K(G,1)$ be a connected CW complex of dimension at most two, and put
\[
 \nu=\max\{\aleph_0,d(H_1(G;\mathbb Z)),
                          \operatorname{rank}H_2(G;\mathbb Z)\}.
\]
For each $g\ne1$ there is a homomorphism $G\to Q$ with
$|Q|\leq\nu^{\aleph_0}$ that does not kill $g$.

\emph{Proof.}
Start with a connected subcomplex of at most $\nu$ cells containing
a loop for $g$, cycles generating $H_1(X)$, and cycles giving a basis
of $H_2(X)$. This last group is free because it is a subgroup of
$C_2(X;\mathbb Z)$. At each of countably many stages, add finite
supports of nullhomotopies for every element of the kernel of the
current fundamental-group map to $G$, and finite bounding chains for
every element of the kernel of its first-homology map to $H_1(X)$.
There are at most $\nu$ choices at each stage. The union $Y$ has at
most $\nu$ cells, and the actual inclusion $Y\subseteq X$ induces
an injection on fundamental groups and an integral homology
isomorphism. Injectivity on $H_2$ requires no choice: there are no
three-boundaries in either complex.

Put $H=\pi_1Y$. Its embedding in $G$ identifies the universal cover
of $Y$ with a component of the inverse image of $Y$ in the
contractible universal cover of $X$. This component has no second
homology, since its cellular two-cycles inject into those of the
contractible two-complex. It is simply connected and
two-dimensional, hence contractible by Hurewicz and Whitehead.
Thus $Y$ itself is a classifying model, and its displayed subgroup
$H$ contains the specified nonidentity $g$.

Collapse a maximal tree of $Y^{(1)}$ and an extension of it to a
maximal tree of $X^{(1)}$. The resulting pair has one vertex and
the relative cellular boundary is an actual isomorphism
\[
 \delta:\mathbb Z^{(R)}\xrightarrow{\cong}\mathbb Z^{(S)},
\]
where $R,S$ are the relative two- and one-cells. For each $r\in R$
choose its finite attaching word over $H*F(S)$ and let $W(r)$ contain
\emph{all} its occurring variables, including those of total exponent
zero. For each $s\in S$, let $I(s)$ be the finite support of
$\delta^{-1}(s)$. A pair of subsets is closed when it is closed under
$r\mapsto W(r)$ and $s\mapsto I(s)$. Closing one additional cell
takes at most countably many cells. By well-ordering the cells this
gives a continuous filtration of actual subcomplexes from $Y$ to $X$
in which each successor adds countably many relative cells. On each
successor quotient $\delta$ remains an isomorphism. Van Kampen
therefore gives presentations whose exponent rows in the newly
added variables are independent, every finite subfamily included.

There is a group $E$ of cardinality at most $\nu$ containing $H$
and solving every finite nonsingular system over itself: perform
countably many passes through all finite systems over the group
present at the beginning of a pass, adjoining universal solutions.
Statement E.1 keeps all coefficient maps injective. Only $\nu$
many finite strings and generators occur in a pass.

Use the single reduced power
\[
 U=E^{\mathbb N}/E^{(\mathbb N)},\qquad
 E^{(\mathbb N)}=\{(e_i):e_i=1\text{ for all but finitely many }i\}.
\]
The denominator is normal; constant sequences embed $E$, and
$|U|\leq\nu^{\aleph_0}$. Any countable nonsingular system over $U$
is solvable in $U$: fix representatives of its coefficients, solve
the first $k$ equations in $E$ at coordinate $k$, and assemble the
solutions coordinatewise. Each equation holds at every sufficiently
large coordinate. Only the finitely many variables occurring in the
first $k$ equations have to be solved there, so E.1 applies exactly.

Now extend the displayed embedding $H\hookrightarrow U$ through
the filtration of the relative presentation. At each successor,
apply the current homomorphism to the coefficients and solve the
resulting countable system in $U$. Assigning the new variables to
these solutions gives a homomorphism on the next presentation with
\emph{exactly} the previous restriction. At a limit use the group
colimit universal property. The resulting homomorphism $G\to U$
is injective on $H$, so it detects $g$. Taking its image as $Q$
proves the assertion. No unproved injection of an intermediate
presentation group is used in this recursion.\hfill$\square$

\paragraph{Statement E.3 (confirmation of INDLOCAL A.20).}
Let $W$ be simply connected, two-dimensional, and
$H_2(W;\mathbb Z)\cong\mathbb Z^{(\kappa)}$ with $\kappa>0$.
There is no bound-$\mu$ two-dimensional resolution when
\[
 \mu\geq \max(\aleph_0,\kappa)^{\aleph_0}.
\]
Indeed any such group has $H_1=0$ and $H_2\cong\mathbb Z^{(\kappa)}$,
so is nontrivial. Statement E.2 detects any nonidentity element in
a quotient of the forbidden size. For $W=S^2$ the excluded range
is exactly $\mu\geq\mathfrak c$ from this argument; the argument
does not exclude $\aleph_0\leq\mu<\mathfrak c$.

\paragraph{Statement E.4 (an explicit finite-bound resolution).}
Every simply connected two-dimensional CW complex has a
two-dimensional resolution satisfying every finite bound at once.
For a target with finite rank $k>0$ in degree two, such a resolution
has a finite classifying model. A target of infinite rank $\kappa$
has one with at most $\kappa$ cells.

\emph{Proof.}
We give an explicit Higman-block proof of this case of the
finite-bound construction in \cite[Sections~8--10]{EBaumslagDyerHeller1980}.
Let
\[
 H=\langle a,b,c,d\mid a^{-1}ba=b^2,
 b^{-1}cb=c^2,\ c^{-1}dc=d^2,\ d^{-1}ad=a^2\rangle.
\]
Its presentation complex is aspherical. To check this, first
amalgamate the standard aspherical Baumslag--Solitar complexes on
$(a,b)$ and $(b,c)$ along the circle for $b$. The subgroup
$\langle a,c\rangle$ is free of rank two: nonzero powers of
$a$ and $c$ lie outside the amalgamated cyclic subgroup, so their
alternating words are reduced. Do the same on $(c,d,a)$, then
amalgamate the two resulting spaces along that free rank-two
subgroup. Injective amalgam normal forms identify the universal
cover as a tree of contractible spaces. The displayed four-relator
presentation is the associated two-dimensional model. Its ordinary
boundary $C_2\to C_1$ is a signed permutation matrix, so it is
acyclic, and each named generator has infinite order.

Every homomorphism $H\to F$ to a finite group is trivial. If any
generator image were nonidentity, choose the least prime $p$
dividing one of their orders, say the order of $x$. The preceding
image $y$ conjugates $x$ to $x^2$. Thus $x$ has odd order, and
the multiplicative order of $2$ modulo $p$ divides the order of
$y$. It is greater than one and divides $p-1$, contradicting
minimality of $p$.

For a cardinal $\kappa>0$, take one central copy of $H$ and
$\kappa$ leaf copies of $H$, amalgamating the infinite cyclic
subgroups generated by $a$ along the edges of a star. Denote
the resulting group by $G_\kappa$. Its graph-of-spaces model
uses only two-dimensional vertex spaces and cylinders on circles,
and is aspherical by the same normal-form argument. The
Mayer--Vietoris sequence for this tree gives
\[
 H_1(G_\kappa;\mathbb Z)=0,\qquad
 H_2(G_\kappa;\mathbb Z)\cong\mathbb Z^{(\kappa)},\qquad
 H_i(G_\kappa;\mathbb Z)=0\quad(i>2).
\]
Every finite image kills every vertex group, and hence all of
$G_\kappa$. The finite boundary matrices for the first eight stars
are independently checked in \texttt{code/E\_higman\_homology.py}.
This calculation is a homology check; asphericity is established
by the preceding geometric argument.

There is an actual homology equivalence to
$W_\kappa=\bigvee_\kappa S^2$. Collapse a maximal tree to make
the source one-vertex. Since $H_1=0$, its boundary
$C_2\to C_1$ is onto and splits over $\mathbb Z$. The splitting
gives a projection $C_2\to H_2$. Collapse the one-skeleton and
map each resulting two-sphere by the finite pinch-and-degree map
specified by its projected basis vector. This defines a continuous
cellular map to $W_\kappa$ and induces the chosen identity on
$H_2$. It is a homology isomorphism in every degree.
Finally every simply connected two-dimensional $W$ is homotopy
equivalent to $W_\kappa$: choose spherical representatives of a
basis of $H_2(W)$ using Hurewicz, form their wedge map, and apply
the simply connected homology Whitehead theorem. If $\kappa=0$,
$W$ is contractible and the trivial group suffices.\hfill$\square$

\paragraph{Statement E.5 (confirmation of INDLOCAL A.13).}
For every infinite cardinal $\lambda$ there is a simple group
$S_\lambda$ of size $\lambda$ with an actual two-dimensional
classifying space and a homology equivalence
$BS_\lambda\to\bigvee_\lambda S^2$. Consequently it is a
bound-$\mu$ resolution for every $\mu<\lambda$.

\emph{Proof.}
Start with $F_\lambda$. At each pass process every ordered pair
$g\ne1,h$ in the group at the beginning of the pass by adjoining
a free letter $t$ and an HNN letter $s$ satisfying
\[
 s^{-1}ws=wh,\qquad
 w=gtg^{-1}t^{-1}gt^2g^{-1}t^{-2}.
\]
In the free product with $\langle t\rangle$, $w$ is cyclically
reduced of syllable length eight. Cyclic reduction of $wh$
has length eight if $hg\ne1$ and six otherwise. Hence both
associated cyclic subgroups are infinite, and HNN normal form
embeds the old group. The relation puts $h$ in the normal
closure of $g$. Repeat the passes countably. Every pair of
elements occurs at a common pass, so the union is simple.

Perform the attachments on classifying spaces. A free letter adds
a circle; each HNN extension adds a circle and a two-cell.
These actual inclusions are aspherical graph-of-spaces extensions,
and their union is aspherical because every sphere map has finite
cell support. The group contains $F_\lambda$ and has at most
$\lambda$ generators, giving cardinality $\lambda$.

In the first pass there are $\lambda$ pairs $(g,1)$. Each of
their distinct two-cells has zero ordinary boundary. They remain
independent in $H_2$ because there are no three-cells. Simplicity
gives $H_1=0$; $H_2$ is free of rank exactly $\lambda$, and all
higher homology vanishes. The split-boundary construction in E.4
now gives the claimed actual map to the wedge. A nontrivial map
from a simple group is injective, which excludes targets of
size less than $\lambda$.\hfill$\square$

\paragraph{Audit E.6 (INDLOCAL A.10 and A.16).}
The general A.10 assertion uses the basic groupoid realization
of Baumslag--Dyer--Heller, Sections 8--10
\cite{EBaumslagDyerHeller1980}. Their Lemma 10.1 is
exactly the reduced-word injectivity needed for its boundary
inclusions, Theorem 9.1 converts the groupoid pushout to the
homotopy pushout, Lemma 10.2 gives the acyclic simplex pieces,
and Theorem 10.4 preserves the dimension for finite simplices.
After killing the inserted Higman factors, each pointed
simplex groupoid is indiscrete. All pointed intersections
contain the common base object; thus gluing those quotients
creates precisely the fundamental groupoid of the original
complex, with no extra loops. The kernel is consequently
normally generated by the inserted Higman groups, and all
finite-group maps kill it. Inductively attaching mapping
cylinders from the already constructed boundary models to
the finite simplex models preserves the required dimension
and cell bound. The local-coefficient comparison is obtained
simplex by simplex and then by cellular unions. These are
the steps needed to supplement the cited dimension theorem;
none is an assertion that arbitrary bar models preserve
dimension. E.4 provides a separate direct verification of
the simply connected two-dimensional cases needed here.

For A.16, the geometric input in A.14 also checks
\cite[Section~6]{EBartholdiMikhailov2026}. Let $L$
be a finite flag triangulation of a positive-genus surface
and form its real moment-angle cube complex $C(L)$. Remove
one open corner to get $M$. Its cut link is that surface.
The dual graph of the three-cells is connected; collapsing
along a spanning tree rooted at a cut face removes every
three-cell. The oriented cube sum maps to a generator in
$H_3(C(L),M)$, so the cut surface maps to zero in $H_2(M)$.
The radial retraction in the CAT(0) universal cover of
$C(L)$ proves injectivity of the old group when $M$ is
glued, through a surface cylinder, to an aspherical space.
After collapsing the lifted old spaces, the universal cover
has flag vertex links (covers of $L$ at the new vertices),
so the quotient is CAT(0). The relative homology sequence
and contractibility of the old lifted pieces make the
uncollapsed universal cover contractible too. The argument
does not require the old aspherical space to have dimension
two, which verifies the extension of the source lemma used
in A.16.

The algebraic steps of A.16 are likewise valid. In the
pullback cover over $\widetilde W$, let
$R=\mathbb Z[\pi_1W]$ and keep the initial copy
$B=H_2(\widetilde W)$. The simple-kernel extension adds
only relative cells in degrees one and two. Its surjective
relative boundary has projective kernel $Z_2$. If the old
$H_1$ has projective dimension at most one, then
$\ker(Z_2\to H_1)$ is projective, by the long exact
$\operatorname{Ext}$ sequence. Hence the new kernel of
$H_2\to B$ is projective. The deliberately included
$\lambda$ commutator relators with $h=1$ split off an
actual free $R^{(\lambda)}$ summand: their coordinates
occur in neither $\partial_2$ nor the old
$\operatorname{im}\partial_3$. The countable direct-sum
swindle makes the whole projective kernel free.

Attach a surface block for each of its basis vectors.
The map of that surface into $\widetilde W$ is nullhomotopic,
because it has zero $H_2$ class and the target is simply
connected. Thus the map to $W$ extends literally. Equivariant
Mayer--Vietoris kills precisely the chosen free kernel,
keeps homology above degree two unchanged, and leaves
$H_1$ of projective dimension at most one and the new
kernel of $H_2\to B$ free. Countable alternation gives
the desired simple kernel and homology equivalence on
pullback covers. The mapping cone of the cellular chain
comparison is an acyclic bounded-below complex of free
$R$-modules, so is contractible; tensoring proves the
comparison for every local system. The dimension bound
is $\max(\dim W,3)$ because the surface cylinders are
three-dimensional. All these steps confirm A.16 and
also identify why it does not answer E2.

\paragraph{Audit E.7 (INDLOCAL A.27--A.31).}
Here $[a,b]=aba^{-1}b^{-1}$. The following records the
individual checks, including the comparisons that preserve
two-dimensional classifying spaces.

For A.27's mixed-identity-free case, the commutator maps
$t\mapsto[a,t]b$ carry a word outside the coefficient
free factor to another word outside it, by malnormality
of a free factor in a free product. Their finite
composition minus the prescribed root value is therefore
a nonidentity mixed word and has a nonidentity evaluation.
This verifies the recursive coefficient choice in
(A.27.4). At a solution whose root equals its $k$th
enumerated value, the $k$th equation becomes constant
and backward substitution contradicts exactly that choice.
For a torsion-free two-conjugacy-class group,
$t\mapsto[a,t]b$ has image $E-\{ab\}$ with infinite
fibres; deleting finitely many inputs preserves its
image. This independently justifies the same recursive
choice in that case. For a uniformly simple group,
the word $W_a=\prod_{i=1}^N v_i[a,u_i]v_i^{-1}$ is
onto whenever $a\ne1$. Indeed a nontrivial commutator
$c=[a,u]$ normally expresses every element in at most
$N$ conjugates, and
$c^{-1}=u[a,u^{-1}]u^{-1}$. The finite tree
substitutions in A.27 are therefore onto, and the
chosen coefficients exclude each candidate root value.

Every displayed relator has exponent one in its own
variable and zero in all others. After lifting its
coefficients from the specified subgroup, finite-prefix
elimination of nonroot variables leaves one unimodular
root relator over the old torsion-free group free-product
with the remaining free letters. Forester--Rourke's
one-pair theorem \cite[Main Theorem]{EForesterRourkeDiagrams}
applies with exactly those hypotheses (an exponent-$\pm1$
shape is amenable):
the coefficient group embeds, and the inclusion-induced
extension-of-scalars map on $\pi_2$ is an isomorphism.
The prefix is therefore aspherical. Compactness of sphere
maps makes the full CW union aspherical; compactness
of nullhomotopies verifies injection of the old group
without needing injections between every two prefixes.
The identity relative boundary gives the integral
homology isomorphism. Killing the coefficient group
first kills the root and then all variables, proving
normal generation. A map extending the chosen subgroup
epimorphism would solve the forbidden system. This
checks the subgroup version required by A.28.

For A.28, a group of size at most $\mathfrak c$ has at
most $\mathfrak c^{\aleph_0}=\mathfrak c$ countable
subgroups. There are at most $\mathfrak c$ countable
group structures and maps, so a pass has at most
$\mathfrak c$ blocks. The $\omega_1$ outer passes
retain that bound. Every map onto a countable group
is onto on a countable subgroup lying in one earlier
pass, since $\operatorname{cf}(\omega_1)>\aleph_0$.
Its listed block contradicts extension of the final
map. This is the precise diagonalization; it never
assumes that a subgroup map extends to the whole
current group. The extra one-variable unimodular
adjunctions preserve asphericity and homology by the
same one-pair theorem, and every finite coefficient
list is eventually processed. Their solvability
passes to quotients simply by lifting coefficients.
Normal generation is transitive through the
attachments and unions.

The Higman image exclusions used in A.28 are valid.
The least-prime argument in E.4 works if the four
generator images merely have finite order, so kills
maps into any torsion group. In a linear image,
each generator is conjugate to its square, forcing
its eigenvalues to be roots of unity. In positive
characteristic it then has finite order. In
characteristic zero, after taking unipotent powers,
the identity
$\operatorname{Ad}(y^{-m})\log(x^L)=2^m\log(x^L)$
forces the logarithm to vanish: conjugation by a
unipotent matrix is unipotent. Thus every generator
again has finite order. A locally linear target
contains the finitely generated image in a linear
subgroup, so it too is killed. Normal generation
transfers these exclusions to the final group.
Mapping each new loop constantly to $S^2$ and
filling each attaching loop extends the original
map literally; the homology isomorphism follows
from the inclusion comparison.

For A.29, nested commutators combine a finite
equational cover of $E-\{1\}$ into a zero-exponent
word $v$ nonidentity at $1$ and identity elsewhere.
Then $v(x)x=1$ is an unsolvable unimodular equation.
Its finite-avoidance consequence supplies the
recursive weighted letters of the proposed metric.
At each stage there are finitely many constraints,
because their conjugators lie in finite balls of
the auxiliary proper length and their word lengths
are bounded. The stage invariant gives
$L(f_j)\geq2^{-j}>0$, while the chosen nonidentity
$a_i$ satisfy $L(a_i)\leq2^{-i}$. Conjugation
changes costs by at most the displayed finite
constant $D_g$, which proves continuity of the
group operations. Thus the claimed nondiscrete
Hausdorff metrizable topology follows. No
completeness or conjugation invariance is obtained.

For A.30, the conjugation-invariant metric does
give the uniform estimate
$\ell(v[b,u]v^{-1})\leq2\ell(b)$. A word with
$L$ variable occurrences is $L$-Lipschitz in the
maximum product metric. Since those conjugated
commutators form a symmetric normal generating
set, a finite path between two different word
values has a first changing step. Choosing $b$
small excludes the prescribed value throughout
the one-step image while retaining two distinct
values. Iteration on the binary tree gives nested
nonempty images excluding every group element.
For each finite prefix, eliminating parents
before children gives an actual homotopy
equivalence relative to the old space with its
wedge of frontier circles. This verifies its
asphericity and coefficient injection independently
of any simultaneous-equation asphericity conjecture.

A.31 replaces the metric estimate by the exact
word-norm bound
\[
 \rho_b\bigl(R(\mathbf s\mathbf u)R(\mathbf u)^{-1}\bigr)
 \leq2L,\qquad
 s_i\in\{v[b,u]v^{-1}\}.
\]
Replacing occurrences one at a time and multiplying
the consecutive differences proves the bound,
including inverse occurrences. Its hypothesis
allows $b$ for which the excluded difference
has greater norm, so the same binary-tree
construction works. If that hypothesis fails,
some $d\ne1$ has $\sup_b\rho_b(d)=M<\infty$.
Expressing any $g$ by $k$ conjugates of $d^{\pm1}$
gives $N_E(g)=\sup_b\rho_b(g)\leq kM$.
Moreover, for every fixed $d\ne1$,
\[
 \rho_d(g)\leq N_E(g)\leq N_E(d)\rho_d(g).
\]
Hence boundedness of any one $\rho_d$ would give
uniform simplicity. For the surviving quotients
all these norms are unbounded; pointwise finiteness
of $N_E$ is nevertheless consistent with this
calculation. A positive invariant length satisfies
$\ell(d)\leq N_E(d)\ell(b)$ for every $b\ne1$,
which proves the claimed discreteness. These
checks confirm the scope of A.31; they do not
eliminate its residual class.

\paragraph{Statement E.8 (the ranges established for E2--E4).}
For $\kappa=0$ every bound is possible. For finite
$\kappa>0$, E.3 and E.4 establish existence for
all finite bounds and nonexistence for
$\mu\geq\mathfrak c$. For infinite $\kappa$,
E.3 and E.5 establish existence for $\mu<\kappa$
and nonexistence for $\mu\geq\kappa^{\aleph_0}$.
In particular if $\kappa^{\aleph_0}=\kappa$,
the exact answer is $\mu<\kappa$.
The outstanding intervals in this work are
\[
 \aleph_0\leq\mu<\mathfrak c
 \quad(0<\kappa<\aleph_0),\qquad
 \kappa\leq\mu<\kappa^{\aleph_0}
 \quad(\kappa\text{ infinite}).
\]
No claim that these intervals are independent of
ZFC is made. Under CH the first interval contains
only $\aleph_0$; this arithmetic observation does
not decide its existence question. The smallest
unresolved case remains $W=S^2$, $\mu=\aleph_0$.

\paragraph{Statement E.9 (limitation of two-dimensional group tests).}
For any cardinal bound, groups with two-dimensional
classifying spaces cannot replace all the group
tests of INDLOCAL B.7 for arbitrary homotopy types.
Indeed, for every actual two-dimensional CW
complex $X$ and every $2$-connected space $Y$,
obstruction theory (or cellular approximation
for $Y=S^3$) gives $[X,Y]=\{*\}$.
In particular $[BG,S^3]$ is a singleton for every
such group, although $S^3$ is not a $1$-type.
Thus E2--E4 can address only the $S^2$ or
$\pi_2$ portion of a bounded-test refinement.
If the test category is the one-object group
$G$, its number of arrows is $|G|$. A nontrivial
bound-$\mu$ resolution usable among tests with
fewer than $\kappa$ arrows must satisfy
$\mu<|G|<\kappa$: the first inequality follows
from the identity homomorphism of $G$. The
missing existence assertion in E2 cannot be
inferred from B.7 or from this necessary size
condition.

\paragraph{Statement E.10 (a discarded proposed construction is impossible).}
There is no nontrivial simple locally free group.
More generally, a nontrivial locally free group
with a finite normal generating set cannot be
perfect. To prove this, choose nonidentity normal
generators $x_1,\ldots,x_r$. If the group were
perfect, write them as finite products of
commutators, then write the finitely many entries
of those commutators as products of conjugates
of the $x_i^{\pm1}$. All the finitely many
elements and conjugators lie in one finitely
generated free subgroup $F$. The Magnus
embedding of $F$ into noncommutative integral
power series assigns every nonidentity element
a least positive degree. Let $d$ be the least
such degree among the $x_i$. Conjugation and
products preserve degree at least $d$, while
the commutator of two such series has degree
at least $2d$. The chosen commutator expressions
therefore give degree at least $2d$ for every
$x_i$, contradicting minimality of $d$.
A nonabelian simple group is perfect and is
normally generated by any nonidentity element,
so the first assertion follows.

\paragraph{Statement E.11 (a restriction on cyclic HNN compression).}
Let $X$ be an actual two-dimensional $K(G,1)$ and
let $G\twoheadrightarrow\mathbb Z$ have kernel $N$.
Put $R=\mathbb Q[t,t^{-1}]$, with $t$ acting by a
generator of the deck group of the cyclic cover
$X_N\to X$. Then $H_2(N;\mathbb Q)$ is a free
$R$-module. If its $R$-rank is $\beta$, there is an
exact sequence
\[
 0\longrightarrow\mathbb Q^{(\beta)}
 \longrightarrow H_2(G;\mathbb Q)
 \longrightarrow
 \ker\bigl(t-1:H_1(N;\mathbb Q)\to H_1(N;\mathbb Q)\bigr)
 \longrightarrow0.
\]
Consequently $\beta\leq\dim_{\mathbb Q}H_2(G;\mathbb Q)$;
if $N$ is perfect, equality holds. Moreover,
\[
 \dim_{\mathbb Q}H_2(N;\mathbb Q)=
 \begin{cases}
 0,&\beta=0,\\
 \max(\aleph_0,\beta),&\beta>0.
 \end{cases}
\]
In particular, a perfect group with uncountable
second Betti number cannot be the cyclic normal
kernel of a two-dimensional group with finite
second Betti number.

For the proof, cellular chains of $X_N$ are free
$R$-modules, including when the sets of cells
are infinite. Since $X_N$ is two-dimensional,
$H_2(X_N;\mathbb Q)=\ker d_2$ is a submodule
of the free module $C_2(X_N;\mathbb Q)$.
The Laurent polynomial ring $R$ is a principal
ideal domain, so this submodule is free, with
some basis of cardinality $\beta$. The short
exact sequence of cellular chain complexes
\[
 0\longrightarrow C_*(X_N;\mathbb Q)
 \xrightarrow{\ t-1\ }C_*(X_N;\mathbb Q)
 \longrightarrow C_*(X;\mathbb Q)\longrightarrow0
\]
gives the displayed exact sequence in homology,
because the coinvariants of a free $R$-module
$R^{(\beta)}$ are $\mathbb Q^{(\beta)}$.
The underlying rational vector space of $R$ has
countably infinite dimension, proving the rank
formula. Finally $X_N$ is a $K(N,1)$, so these
are the group homology groups claimed.

Applied to $G=S\rtimes_f\mathbb Z$, this rules
out reducing an uncountable $H_2(S;\mathbb Q)$
to finite rank merely by choosing an automorphism
$f$ with small coinvariants while retaining an
actual two-dimensional model. For a proper
ascending HNN extension the cyclic kernel is
instead $\operatorname*{colim}(S\xrightarrow fS
\xrightarrow f\cdots)$; the same statement applies
to that kernel. Killing its homology first does
not produce a two-dimensional model for the
extension. If $S$ was simple, that direct limit
is simple, so an acyclic two-dimensional cyclic
kernel would already supply the difficult
uncountable acyclic simple group, rather than
deducing one from a homology calculation.

\paragraph{Statement E.12 (ordinary countable equation compactness fails).}
Let $E$ be a countably infinite group such that, for every $d\in E$,
there is a finite integer $n(d)$ for which $d$ is a product of at most
$n(d)$ conjugates of $a^{\pm1}$ for every $a\ne1$ in $E$.
Then there is a countable system of equations with coefficients in $E$
whose every finite subsystem has a solution in $E$, but which has no
simultaneous solution in $E$.
In particular the pointwise finite alternative for $N_E$ in
INDLOCAL A.31 cannot occur in a countably equation-compact group.

\emph{Proof.}
Fix a nonprincipal ultrafilter $\mathcal U$ on $\mathbb N$ and put
$P=E^{\mathbb N}/\mathcal U$, identifying $E$ with its diagonal copy.
For any $b\ne1$ in $P$ and any $d\in E$, the bound $n(d)$ implies
$d\in\langle\!\langle b\rangle\!\rangle_P$.
Indeed, choose representatives $b_i$. On a set in $\mathcal U$,
$b_i\ne1$, and there are expressions for $d$ as products of at most
$n(d)$ conjugates of $b_i^{\pm1}$.
There are only finitely many possible lengths and sign patterns;
one pattern occurs on a set in $\mathcal U$.
The conjugating elements in that pattern, chosen coordinatewise and
defined arbitrarily off this set, give the required expression in $P$.
Thus every nontrivial normal subgroup of $P$ contains the diagonal $E$.

Choose $x\in P\setminus E$, for instance the class of a sequence
of distinct elements of $E$. Construct countable subgroups
$P_0\le P_1\le\cdots$ of $P$, starting with $P_0=\langle E,x\rangle$.
At step $j+1$, for each $b\in P_j\setminus\{1\}$ and $d\in E$,
adjoin the finitely many conjugators in one expression witnessing
$d\in\langle\!\langle b\rangle\!\rangle_P$.
There are only countably many pairs, so
$P'=\bigcup_jP_j$ is countable and contains $E$ properly.
Every nontrivial normal subgroup of $P'$ contains $E$.

Use one variable for each element of $P'\setminus E$; for an element
of $E$ use that element itself as a coefficient.
Impose all the multiplication equations of $P'$.
This is a countable system. Every finite subsystem has a solution
in $P$, namely the displayed elements of $P'$, and consequently in
$E$: for the finitely many equations, take representatives of these
elements and choose one coordinate in the intersection of the
finitely many sets in $\mathcal U$ on which the equations hold.
If the entire system had a solution in $E$, its values would define
a homomorphism $r:P'\to E$ restricting to the identity on $E$.
Its kernel is nontrivial since
$x r(x)^{-1}\ne1$, whereas every nontrivial normal subgroup of $P'$
contains $E$. This contradicts $r|_E=\mathrm{id}_E$.
\hfill$\square$

The exponent-sum rows of this multiplication-table system are not
asserted independent. The argument supplies neither a nonsingular
system nor an actual two-dimensional relative homology equivalence.
Adding fresh variables $y_i=r_i$ to make arbitrary equations
formally independent removes their restrictions; imposing $y_i=1$
restores the dependency. Thus E.12 gives a precise compactness
obstruction, but does not decide E2.

\paragraph{Unresolved step and approaches that do not close it.}
Replacing the reduced power in E.2 by a countable
universal target requires solutions to countable
systems, not merely finite ones; A.27 rules out
the natural countable existentially closed
targets. Repeating the metric argument of
A.30 with the topology from A.29 loses the
uniform conjugation bound. Pointwise finite
$N_E$ in A.31 permits uniform bounds for
reaching any specified finite set, but not
surjectivity onto the whole target; nested
tree preimages need not admit a common bound.
Finally, replacing old two-cells by the
Bartholdi--Mikhailov blocks can retain
two-dimensional models at each finite step,
but those structural maps are not subcomplex
inclusions: the old cell is removed by a
collapse with a three-cell. A mapping telescope
then has dimension three. These are exact
obstructions to the attempted arguments,
not proofs of nonexistence in E2.
